\chapter{Conclusions}
\label{cha:conclusions}

This work presented an analysis of Linux kernel tracing, with the objective of
understanding the systems and interfaces involved in tracing kernel events, as
well as the evolution of the mechanisms used to implement them. The study
examined the main tracing frameworks available in Linux, with particular
attention to ftrace and perf, and explored the events exposed by the kernel and
the capabilities offered by these tools. It also considered the role of modern
tracing applications, focusing on eBPF. The analysis provided an overview of the
underlying mechanisms and their implementation, examining the relevant source
code and the technical details necessary to understand their behaviour.

Linux tracing remains a continuously evolving field, reflecting the ongoing
development of the kernel and the increasing demands placed on observability,
performance analysis, and security. New features, changes to existing
interfaces, and improvements to the underlying implementation require tracing
tools and the knowledge surrounding them to evolve accordingly. For this reason,
the results presented in this work should be considered within the context of a
system that is continually changing, rather than as a definitive description of
a fixed set of mechanisms.

It is hoped that this analysis will prove useful to both developers already
working with Linux tracing and those approaching the subject for the first time,
helping them navigate the relevant interfaces, understand the underlying
mechanisms, and make informed decisions when developing or maintaining tracing
tools. In this respect, the work aims to contribute to the knowledge required to
follow the evolution of Linux tracing and to support further development in this
area.
